% Archiv: Aufgaben aus "Von konkreten Problemen zur Wahrscheinlichkeitstheorie" (ehem. S1.2)
% Nicht in aktueller Dateistruktur eingebunden -- zur späteren Verwendung.

% ===========================================================================
% AUFGABEN
% ===========================================================================

% ---------------------------------------------------------------------------
\subsection*{Diskrete vs. Stetige Zufallsexperimente}

\begin{komplexaufgaben}{sec:historischeProblem}

	\komplexaufgabe{
		\textbf{Diskrete vs. Stetige Zufallsexperimente}

		Betrachten Sie folgende Situationen:
		\begin{enumerate}[label=(\Roman*)]
			\item Würfeln mit einem fairen Würfel
			\item Wartezeit an einer Ampel (0 bis 60 Sekunden)
			\item Anzahl der Anrufe in einer Telefonzentrale pro Stunde
			\item Punkt zufällig in einem Quadrat wählen
			\item Anzahl defekter Glühbirnen in einer Packung mit 10 Stück
			\item Lebensdauer einer Glühbirne (in Stunden)
			\item Temperatur um 12 Uhr mittags (gemessen in $^\circ$C)
		\end{enumerate}

		\textbf{Arbeitsaufträge:}
	}{
		\item \textbf{Kategorisierung:} Ordnen Sie die Situationen (I)-(VII) ein:
		\begin{itemize}
			\item Diskret oder stetig?
			\item Falls diskret: endlich oder unendlich viele Ausgänge?
		\end{itemize}

		\item \textbf{Ergebnismengen:} Geben Sie für jede Situation die Ergebnismenge $\Omega$ an.

		\item \textbf{Gleichwahrscheinlichkeit:} Bei welchen Situationen ist ein Laplace-Modell sinnvoll?

		\item \textbf{Gemeinsame Fragen:} Formulieren Sie für drei Situationen je eine Frage der Form „Mit welcher Wahrscheinlichkeit...?" Was haben diese Fragen gemeinsam? Was unterscheidet sie?

		\item \textbf{Mathematische Herausforderung:} Welche Situationen können Sie mit Ihrem bisherigen Wissen bearbeiten? Bei welchen stoßen Sie an Grenzen?

		\item \textbf{Bedarf an Verallgemeinerung:} Warum braucht man eine einheitliche mathematische Theorie?
	}

\end{komplexaufgaben}

% ---------------------------------------------------------------------------
\subsection*{Das Buffonsche Nadelproblem}

\begin{komplexaufgaben}{sec:historischeProblem}

	\komplexaufgabe{
		\textbf{Das Buffonsche Nadelproblem (Experiment und Theorie)}

		Auf einem Blatt Papier sind parallele Linien im Abstand $d$ gezeichnet. Eine Nadel der Länge $\ell < d$ wird zufällig auf das Blatt geworfen.

		\textbf{Arbeitsaufträge:}
	}{
		\item \textbf{Experiment:} Führen Sie das Experiment durch (oder simulieren Sie es):
		\begin{itemize}
			\item Zeichnen Sie parallele Linien im Abstand $d = 3$ cm.
			\item Verwenden Sie eine „Nadel" (z.B. Zahnstocher) der Länge $\ell = 2$ cm.
			\item Werfen Sie die Nadel 50-mal und notieren Sie, wie oft sie eine Linie schneidet.
		\end{itemize}

		\item \textbf{Modellierung:} Beschreiben Sie die Position der Nadel durch:
		\begin{itemize}
			\item Den Abstand $x$ des Nadelmittelpunkts zur nächsten Linie ($0 \le x \le \frac{d}{2}$)
			\item Den Winkel $\alpha$ der Nadel zu den Linien ($0 \le \alpha \le \frac{\pi}{2}$)
		\end{itemize}

		\item \textbf{Schnittbedingung:} Die Nadel schneidet eine Linie, wenn $x \le \frac{\ell}{2} \sin(\alpha)$. Machen Sie sich dies geometrisch klar.

		\item \textbf{Theoretische Wahrscheinlichkeit:}
		\[
		P(\text{Schnitt}) = \frac{\int_0^{\pi/2} \frac{\ell}{2}\sin(\alpha)\,d\alpha}{\frac{d}{2} \cdot \frac{\pi}{2}}
		\]
		Berechnen Sie dieses Integral und zeigen Sie: $P(\text{Schnitt}) = \dfrac{2\ell}{\pi d}$

		\item \textbf{Überraschung:} Für $\ell = d$ gilt $P(\text{Schnitt}) = \frac{2}{\pi}$. Wie kann man damit $\pi$ schätzen?

		\item \textbf{Unterschied zu Würfeln:} Warum funktioniert hier „Abzählen" nicht? Welche neue Methode war nötig?
	}

\end{komplexaufgaben}

% ---------------------------------------------------------------------------
\subsection*{Reflexion: Verschiedene Modelle, gemeinsame Struktur}

\begin{komplexaufgaben}{sec:historischeProblem}

	\komplexaufgabe{
		\textbf{Reflexion: Verschiedene Modelle, gemeinsame Struktur}

		Sie haben nun verschiedene Arten von Wahrscheinlichkeitsproblemen kennengelernt:
		\begin{itemize}
			\item Diskret mit endlich vielen Ausgängen (Würfel, Münze)
			\item Diskret mit unendlich vielen Ausgängen (Anzahl Anrufe)
			\item Stetig/geometrisch (Buffon, Wartezeiten)
		\end{itemize}

		\textbf{Arbeitsaufträge:}
	}{
		\item \textbf{Gemeinsamkeiten:} Was haben alle diese Situationen gemeinsam? (z.B. „Wahrscheinlichkeiten liegen zwischen 0 und 1")

		\item \textbf{Unterschiede:} Erstellen Sie eine Tabelle für die drei Kategorien (Beispiel, Methode, Werkzeuge).

		\item \textbf{Rechenregeln:} Verifizieren Sie für je ein Beispiel aus jeder Kategorie:
		\begin{itemize}
			\item $P(A \cup B) = P(A) + P(B)$, falls $A$ und $B$ unvereinbar sind
			\item $P(A^c) = 1 - P(A)$
		\end{itemize}

		\item \textbf{Bedarf an Abstraktion:} Welche Vorteile hätte ein einheitlicher, abstrakter Rahmen?

		\item \textbf{Ausblick auf Axiome:} Was erwarten Sie von den Axiomen von Kolmogorov?
	}

\end{komplexaufgaben}

% ---------------------------------------------------------------------------
\subsection*{Enrichment: Das Bertrand-Paradoxon}

\begin{komplexaufgaben}{sec:historischeProblem}

	\komplexaufgabe{
		\textbf{Das Bertrand-Paradoxon: Präzision ist alles!}

		Gegeben ist ein Kreis mit Radius $r$. Eine Sehne wird \emph{zufällig} im Kreis gewählt.

		\textbf{Frage:} Mit welcher Wahrscheinlichkeit ist die Sehne länger als die Seite des einbeschriebenen gleichseitigen Dreiecks (Länge $r\sqrt{3}$)?

		\textbf{Arbeitsaufträge:}
	}{
		\item \textbf{Methode 1 – Zufälliger Endpunkt:} Ein Punkt ist fest, der andere zufällig auf dem Kreis. Berechnen Sie die Wahrscheinlichkeit.

		\item \textbf{Methode 2 – Zufälliger Mittelpunkt:} Der Mittelpunkt der Sehne ist zufällig im Kreis. Berechnen Sie die Wahrscheinlichkeit (Kreisflächenvergleich).

		\item \textbf{Methode 3 – Zufälliger Abstand:} Richtung fest, Abstand $d$ zum Mittelpunkt gleichverteilt in $[0, r]$. Berechnen Sie die Wahrscheinlichkeit.

		\item \textbf{Vergleich:} Sie erhalten $\frac{1}{3}$, $\frac{1}{4}$ und $\frac{1}{2}$. Welche Antwort ist „richtig"?

		\item \textbf{Das Paradoxon:} Warum führen drei scheinbar gleich natürliche Methoden zu verschiedenen Ergebnissen?

		\item \textbf{Lehre:} Was lernen wir über die Notwendigkeit präziser Definitionen?

		\item \textbf{Verbindung zu $(\Omega, \mathcal{F}, P)$:} Welcher der drei Teile war bei den drei Methoden unterschiedlich?
	}

\end{komplexaufgaben}


% ===========================================================================
% LÖSUNGEN
% ===========================================================================

\subsection*{Lösungen: Diskrete vs. Stetige Zufallsexperimente}

\begin{komplexloesungen}{sec:historischeProblem}

	\komplexloesung{
		\textbf{Diskrete vs. Stetige Zufallsexperimente}
	}{
		\item (I) diskret-endlich; \enspace (II) stetig; \enspace (III) diskret-unendlich; \enspace
		      (IV) stetig; \enspace (V) diskret-endlich; \enspace (VI) stetig; \enspace (VII) stetig

		\item (I) $\Omega = \{1,2,3,4,5,6\}$; \enspace
		      (II) $\Omega = [0,60]$; \enspace
		      (III) $\Omega = \{0,1,2,\ldots\}$; \enspace
		      (IV) Einheitsquadrat $[0,1]^2$; \enspace
		      (V) $\Omega = \{0,1,\ldots,10\}$; \enspace
		      (VI) $\Omega = [0,\infty)$; \enspace
		      (VII) $\Omega = \mathbb{R}$

		\item Laplace sinnvoll: (I) fairer Würfel; (V) falls alle Birnen gleich hergestellt werden.

		\item \emph{(Offene Lernendendiskussion)}

		\item Mit Kombinatorik/Laplace bearbeitbar: (I), (V). Grenzen bei stetigen und unbeschränkten Modellen.

		\item Eine einheitliche Theorie ermöglicht denselben Rechenrahmen für alle Modelltypen.
	}

\end{komplexloesungen}

\subsection*{Lösungen: Das Buffonsche Nadelproblem}

\begin{komplexloesungen}{sec:historischeProblem}

	\komplexloesung{
		\textbf{Das Buffonsche Nadelproblem}
	}{
		\item \emph{(Experimentelles Ergebnis variiert; Erwartungswert: ca.\ 67\,\% Treffer bei $\ell=2$, $d=3$)}

		\item $x$ und $\alpha$ reichen, weil Mittelpunktlage und Winkel die Nadel eindeutig beschreiben.

		\item Schnitt genau dann, wenn $x \le \tfrac{\ell}{2}\sin\alpha$.

		\item $\displaystyle\int_0^{\pi/2} \tfrac{\ell}{2}\sin\alpha\,d\alpha = \tfrac{\ell}{2}$; \quad
		      Fläche des Rechtecks: $\tfrac{d\pi}{4}$; \quad
		      $P = \dfrac{2\ell}{\pi d}$ \checkmark

		\item Für $\ell=d$: $P = \tfrac{2}{\pi}$. Schätzung: $\pi \approx \dfrac{2 \cdot \text{Würfe}}{\text{Treffer}}$.

		\item Würfeln: diskret, Abzählen möglich. Buffon: stetig, Wahrscheinlichkeit = Flächenverhältnis.
	}

\end{komplexloesungen}

\subsection*{Lösungen: Reflexion}

\begin{komplexloesungen}{sec:historischeProblem}

	\komplexloesung{
		\textbf{Reflexion: Verschiedene Modelle, gemeinsame Struktur}
	}{
		\item $P(E) \in [0,1]$; Gesamtwahrscheinlichkeit 1; Komplementregel; Additivität bei unvereinbaren Ereignissen.

		\item \emph{(Tabelle — Lernendenergebnis)}

		\item \emph{(Konkrete Beispielverifikation — Lernendenergebnis)}

		\item Einzeltheorien wären redundant; einheitlicher Rahmen ist effizienter und allgemein übertragbar.

		\item Axiome: Nicht-Negativität ($P \ge 0$), Normierung ($P(\Omega)=1$), Additivität bei unvereinbaren Ereignissen.
	}

\end{komplexloesungen}

\subsection*{Lösungen: Das Bertrand-Paradoxon}

\begin{komplexloesungen}{sec:historischeProblem}

	\komplexloesung{
		\textbf{Das Bertrand-Paradoxon}
	}{
		\item Methode~1: $P = \tfrac{1}{3}$.

		\item Methode~2: $P = \dfrac{\pi(r/2)^2}{\pi r^2} = \dfrac{1}{4}$.

		\item Methode~3: $P = \tfrac{1}{2}$.

		\item Alle drei Rechnungen sind korrekt — das Ergebnis hängt vom gewählten Zufallsmechanismus ab.

		\item Der Begriff „zufällig" ist ohne genaue Festlegung nicht eindeutig.

		\item Präzise Definitionen sind unerlässlich; „wähle zufällig" allein ist keine vollständige Beschreibung.

		\item Unterschiedlich war die Wahrscheinlichkeitsvorschrift $P$.
	}

\end{komplexloesungen}
